% Original: PDF strana 1, donji slajd.
\slajdid{03-s002}
\begin{frame}[label=03-s002]{Stepeni integracije}
% element: 03-s002:e01
Podela integrisanih kola zasniva se na broju komponenti (tranzistora) i logičkih kapija na jednom čipu.
\par\vspace{3mm}
% element: 03-s002:e02
{\centering
\begingroup
\fontsize{10}{12}\selectfont
\setlength{\tabcolsep}{6pt}\renewcommand{\arraystretch}{1.3}
\begin{tabular}{@{}llcc@{}}\toprule
\textbf{Stepen} & \textbf{Naziv} & \textbf{Broj komponenti} & \textbf{Broj kapija} \\ \midrule
\textbf{SSI} & Small Scale Integration & do 100 & do 10 \\
\textbf{MSI} & Medium Scale Integration & do 500 & 10--100 \\
\textbf{LSI} & Large Scale Integration & do 300\,000 & preko 100 do 10\,000 \\
\textbf{VLSI} & Very Large Scale Integration & preko 300\,000 & 10\,000--100\,000 \\
\textbf{ULSI} & Ultra Large Scale Integration & preko 1\,500\,000 & preko 100\,000 \\ \bottomrule
\end{tabular}\par
\endgroup}
\par\vspace{3mm}
% element: 03-s002:e03
\Rezultat{Granice nisu stroge. Pojam VLSI dizajn obuhvata i VLSI i ULSI tehnologije.}
\end{frame}
